% ============================================================
% Chapter 6: Experimental Design
%   Section: Construction of the Custom Validation Set
%     Subsection: Synthetic Refuted Claims
% Included from chapters/06-experimental-design/03-custom-validation-set/custom-validation-set.tex
% ============================================================

% Source: labeller/app.py (MUTATIONS, Store.mutate, stale_mutations),
% backend/src/evaluation/dataset.py (origin), and the 30 records
% fact149-fact178 in backend/data/evaluation/manual/.

\subsection{Synthetic Refuted Claims}
\label{subsec:custom-mutations}

The natural claims contain no refuted one
(Table~\ref{tab:custom-matrix}). Constructive outlets rarely publish a
claim the evidence contradicts, and the claims they do get wrong are
more often wrong in a detail (\texttt{PARTIALLY\_TRUE}) than in their
substance. A verification benchmark without a refuted class cannot
measure whether a system detects false information, so the
\texttt{FALSE} class is built, as the benchmarks this work follows
built theirs.

\paragraph{The precedent.} FEVER obtained most of its refuted claims
by \emph{mutating} claims extracted from Wikipedia: annotators
paraphrased, negated, generalised or specified them, or substituted an
entity or a relation with a similar or a dissimilar one, and the
mutated claims were then labelled against the evidence
\cite{thorne2018fever}. SciFact had expert annotators write the
negation of claims supported by a scientific abstract, so that the same
abstract refutes them \cite{wadden2020scifact}. Both methods share
the property that matters here: the evidence that supports the
original claim is the evidence that refutes the mutated one, so the
refuted label rests on a source already checked, not on a new search.

\paragraph{What was learned from the precedent.} Mutation leaves
traces in the wording. In FEVER, negation words such as ``not'' and
``only'' were frequent enough in refuted claims that a model reading
the claim alone, without any evidence, reached an accuracy close to
that of evidence-based systems \cite{schuster2019debiasing}. A
benchmark with that artefact measures a model's sensitivity to
wording, not its verification. The procedure below avoids it by
design.

\paragraph{The procedure.}
\begin{enumerate}[leftmargin=*]
    \item \textbf{Source.} Only a natural claim labelled
    \texttt{TRUE}, whose note states the value the evidence gives, is
    mutated; each source is mutated once.
    \item \textbf{Substitution only, of what the claim asserts.} One
    element is replaced: a number, a date, an entity or a relation
    (\emph{raised} for \emph{lowered}, the overall prize for the young
    photographer's category). Negation is not used. The replaced
    element must be what the claim asserts, not a detail of it: under
    the guide's numeric rule (Subsection~\ref{subsec:custom-rules}), a
    wrong detail of an otherwise correct claim is
    \texttt{PARTIALLY\_TRUE}, not \texttt{FALSE}. Where the claim
    states figures that must agree with each other, they are changed
    together, so that the claim does not refute itself.
    \item \textbf{Plausible values.} The substitute is chosen to be
    plausible without the evidence --- two million new cases of breast
    cancer a year is roughly the world's figure, not Latin America's;
    the first Lagrange point lies as far from Earth as the second ---
    so that knowledge of the world, or the evidence, and not the form
    of the claim, is what refutes it.
    \item \textbf{Everything else unchanged.} The mutated record keeps
    its source's article, date, topic, claim type, source tier and
    evidence links, and names its source and the kind of substitution
    (\texttt{mutatedFrom}, \texttt{mutation}). The labelling form
    refuses a mutation whose source is not a natural \texttt{TRUE}
    claim, or whose text is unchanged, and lists any mutation whose
    source is later relabelled, since its label was only valid by
    construction from that source.
    \item \textbf{Stratified, and reported apart.} Thirty refuted
    claims were made, six per topic group, nine of them in Spanish;
    twelve substitute a number, eight a date, six an entity and four a
    relation. The evaluation reports every metric for natural and
    synthetic claims separately (the \texttt{origin} breakdown), as
    well as together, because a system can handle substituted claims
    well and the claims the world produces badly.
\end{enumerate}

\paragraph{What the synthetic class can and cannot show.} It shows
whether the system notices when a claim contradicts evidence it can
find, which is the core of verification. It does not show how often
real constructive news is false, nor the kinds of error real articles
make; the natural claims answer those questions, and the X-Fact
benchmark supplies refuted claims the world actually produced
\cite{gupta2021xfact}. Like the AI-assisted labels of
Subsection~\ref{subsec:custom-procedure}, the thirty mutations were
proposed by a language model and are reviewed by the annotator; the
blind re-label (Subsection~\ref{subsec:custom-review}) draws from
them as from any other claim, which also tests whether a person reading
the evidence reaches \texttt{FALSE} where the construction says so.
